\documentclass[10pt,a4paper]{amsart}
\usepackage[margin=19mm]{geometry}
\usepackage{amsmath,amssymb,mathtools}
\usepackage[scr=rsfs,cal=euler]{mathalfa}
\usepackage{xfrac}
\usepackage[unicode,hidelinks,bookmarks=false]{hyperref}
\newcommand{\ie}{i.e.\ }
\newcommand{\cf}{cf.\ }
\newcommand{\resp}{respectively\ }
\newcommand{\Subsection}{\subsection}
\providecommand{\nobreakdash}{\nobreak\hbox{-}}
\setlength{\emergencystretch}{2em}
\multlinegap0pt
\usepackage{amssymb}
\usepackage{xcolor}
\usepackage{tcolorbox}
\newtheorem{lemma}{Lemma}
\newtheorem{claim}{Claim}
\title{Some applications of the real strict order property hierarchy}
\author{Scott Mutchnik}
\date{Source: arXiv:2606.28740v2 (CC BY 4.0)}
\begin{document}
\maketitle

\noindent\emph{Source and licence.} Some applications of the real strict order property hierarchy. Version arXiv:2606.28740v2 (CC BY 4.0), distributed by arXiv under the Creative Commons Attribution 4.0 International licence, http://creativecommons.org/licenses/by/4.0/, as recorded in the arXiv licence metadata for this exact version and shown on the abstract page https://arxiv.org/abs/2606.28740v2. Full licence notice: \texttt{licenses/arxiv-2606.28740-CC-BY-4.0.txt}.\par\smallskip

\noindent\textit{Editorial scope.} Counted unit: the article's lemma tp-structure (source0.tex lines 1512--1514). The article's complete original proof of it is delivered with it (source0.tex lines 1517--1612, one \texttt{proof} environment of 96 lines). No mathematics is written, repaired or re-derived: the statement and the proof are the article's own lines, in the article's own notation, and no retained line is substituted or re-flowed.  No further source-local context is needed: every cross-reference of every retained line resolves inside the two delivered spans.  Nothing was omitted from any retained span, so the package carries no omission record.  No retained line contains a citation, so the wrapper carries no bibliography and no bibliography entry.  No retained line contains a figure environment, an image-inclusion command or drawing code, so no image is delivered and the package declares no asset dependency.  Editorial formatting only, in addition: an AMS \texttt{amsart} preamble that restates the article's own declarations and macros used by the retained lines, namely the environments \texttt{lemma}, \texttt{claim} on separate counters, together with the standard packages those lines need.  The retained environments print in this document as Lemma 1, Claim 1 in order of appearance, whereas the article prints the same items with the numbering of its own full document.  The complete native source of the article as distributed in the arXiv e-print for this exact version is bundled unchanged as \texttt{sources/203876/source0.tex}.
\begin{lemma}
    \label{embedding into the standard tp-structure} Let $A$ be a finite presidium with no alternating cycles. Then $A$ (weakly) embeds into the standard $\mathrm{TP}$-structure.
\end{lemma}

\begin{proof}
    We first prove the following:


        \begin{claim}
       \label{order construction 1} Let $A$ be a presidium with no alternating cycles. Let $a, b$ be distinct elements of $P(A)$ belonging to a common convivium, and let $c, d$ be distinct elements of $P(A)$ belonging to a different common convivium. Suppose that there is $e \in O(A) $ such that $R(e, b)$ and $R(e, c)$. Then either the presidium $A' : = A \sqcup \{f\}$, $f \in O(A')$ whose instances of $R, I$ are exactly those of $A$ together with $R(f, a)$ and $R(f, c)$ has no alternating cycles, or the presidium $A' : = A \sqcup \{f\}$, $f \in O(A')$ whose instances of $R, I$ are exactly those of $A$ together with $R(f, b)$ and $R(f, d)$ has no alternating cycles.
    \end{claim}
\begin{proof}
    (of claim)
    
    Suppose that the first of the two possibilities stated in the claim fails: the presidium $A' : = A \sqcup \{f\}$, $f \in O(A')$ whose instances of $R, I$ are exactly those of $A$ together with $R(f, a)$ and $R(f, c)$ does have an alternating cycle. Then because $A$ itself does not already have any alternating cycles, there is a sequence 
   
   $$a_1, a_2, a^{1}, a_{3}, a_{4}, a^{2}, \ldots a_{2k-1}, a_{2k}, a^{k}, a_{2k+1}, a_{2k+2}, \ldots a_{2n-1}, a_{2n}, a^{n}, a_{2n+1}, a_{2n+2} $$

for $n \geq 1$ consisting of distinct $a_i \in P(A)$, $a^i \in O(A)$ with $I^{*}(a_i, a_{i+1})$ for $i \leq 2n+2$ odd, $R(a^i, a_{2i})$ and $R(a^i, a_{2i+1})$ for $i \leq n$,  and $a_1 = a$ and $a_{2n+2} = c$.  We claim $a_2 = b$. Suppose otherwise, and suppose first that one of the $a_i$ (but not $a_1= a$ or $a_2$) is equal to $b$; say this is $a_k$.  In the case that $k$ is odd, we have $I^{*}(a_k, a_2)$ because $I^{*}(a, b) $ and $I^{*}(a_{1}, a_{2})$ for $a_{1} = a$, so $a_{2}$ and $a_{k} = b$ are distinct members of the same convivium. So the interval of the $a_i, a^{i}$, as ordered according to the above enumeration, between $a_{2}$ and $a_{k}$ inclusive will comprise an alternating cycle, a contradiction.  In the case that $k$ is even, we have $I^{*}(a_k, a_2)$ as before, but also $I^{*}(a_{k-1}, a_k)$, so $a_{2}$ and $a_{k-1}$ belong to the same convivium and $I^{*}(a_{k-1}, a_{2})$. Thus, the interval of the $a_i, a^{i}$ between $a_{2}$ and $a_{k-1}$ inclusive will comprise an alternating cycle, again a contradiction.  So given our supposition that $a_2 \neq b$, we see that none of the $a_i$ are equal to $b$. Now suppose one of the $a^i$, say $a^{k}$, is equal to $e$.  Then as before, $I^{*}(b, a_{2})$, while also $R(a^{k}, b)$. By this and the fact that none of the $a_i$ are equal to $b$, the interval between $a_{2}$ and $a^{k}$ inclusive together with $b$ will comprise an alternating cycle, giving us a contradiction yet again. So because none of the $a_i$ are equal to $b$ or $a_i$ to $e$, while $R(e, a_{2n+2})$ (because $a_{2n+2}=c$), $R(e, b)$ and $I^{*}(b, a_1)$, the $a_i$ and $a^i$ together with $e$ and $b$ will comprise an alternating cycle, another contradiction.

We have shown that $a_2 = b$; before proceeding, we will also show that none of the $a^{i}$ are equal to $e$. Suppose, say, that $a^{k} = e$; then $R(a^{k}, a_{2n+2})$. So we get another contradiction: the interval between $a^k$ and $a_{2n+2}$ inclusive comprises an alternating cycle.

In summary, taking the interval between $a_2$ and $a_{2n+2}$ inclusive in the above enumeration of the $a_i$ and $a^{i}$, we get a sequence


$$h_{2}, h^{1}, h_{3}, h_{4}, h^{2}, \ldots h_{2k-1}, h_{2k},  h^{k}, h_{2k+1}, h_{2k+2}, \ldots h_{2n-1}, h_{2n},  h^{n}, h_{2n+1}, h_{2n+2}$$

for $n \geq 1$ consisting of distinct $h_i \in P(A)$ and $h^i \in O(A)$ with $I^{*}(h_i, h_{i+1})$ for odd $i$ with $2 \leq i \leq  2n+2$, $R(h^{1}, h_2)$ and $R(h^1, h_3)$, $R(h^i, h_{2i})$ and $R(h^i, h_{2i+1})$ for $i = 2, \ldots n $, $h_{2}= b $ and $h_{2n+2} = c$, and all of the $h^{i}$ distinct from $e$.


Now suppose the second of the two possibilities stated in the claim also fails. Then similarly, we get a sequence

$$g_{2}, g^{1}, g_{3}, g_{4}, g^{2}, \ldots g_{2k-1}, g_{2k},  g^{k}, g_{2k+1}, g_{2k+2}, \ldots g_{2m-1}, g_{2m},  g^{m}, g_{2m+1}, g_{2m+2}$$

for $m \geq 1$ consisting of distinct $g_i \in P(A)$ and $g^i \in O(A)$ with $I^{*}(g_i, g_{i+1})$ for odd $i$ with $2 \leq i \leq  2m+2$, $R(g^{1}, g_2)$ and $R(g^1, g_3)$, $R(g^i, g_{2i})$ and $R(g^i, g_{2i+1})$ for $i = 2, \ldots m $, $g_{2}= c $ and $g_{2m+2} = b$, and all of the $h^{i}$ distinct from $e$. (So all that is different about these conditions is that $b$ and $c$ are reversed.)

Now suppose that the intersection of these two sequences is more than just $\{b, c\}$. Suppose first that the first of the $g_i, g^i$ besides $g_2=b$ in this intersection, per the above enumeration, is of the form $g^i$, say $g^k$, and that it coincides with $h^{\ell}$. Then the interval between $g_2$ and $g^k$ inclusive (again per that enumeration), together with the interval between $h^{\ell} = g^k$ and $h_{2n+2} = g_2$ (per the above enumeration on the $h_i, h^i$), will comprise an  alternating cycle, a contradiction. Now suppose that the first of the $g_i, g^i$ besides $g_2$ in the intersection is of the form $g_i$, say $g_k$, and it coincides with $h_\ell$. First consider the cases where $k$ is even and $\ell$ is even, or $k$ is odd and $\ell$ is odd. Then the interval between $g_2$ and $g_k$ inclusive, together with the interval between $h_{\ell} = g_k$ and $h_{2n+2} = g_2$, will comprise an alternating cycle, again a contradiction. Next, consider the case where $k$ is even and $\ell$ is odd. Then $I^{*}(g_{k-1}, g_{k})$ and $I^{*}(h_{\ell}, h_{\ell+1})$, so $g_{k-1}$ and $h_{\ell+1}$ share a convivium and $I^{*}(g_{k-1}, h_{\ell+1})$. Then the interval between $g_2$ and $g_{k-1}$ inclusive, together with the interval between $h_{\ell+1}$ inclusive and  $h_{2n+2} = g_2$, will comprise an alternating cycle, again a contradiction. Lastly, consider the case where $k$ is odd and $\ell$ is even. Then, using that all of the $g^i$, $h^i$ are distinct from $e$, the interval between $g_{2}=c$ and $g_{k}$, inclusive, together with the interval between $h_{2}=b$ inclusive and $h_{\ell} = g_{k}$ (which will be read in reverse), together with $e$, will comprise an alternating cycle, yet again a contradiction.

So the intersection of the sequence of $g_i, g^i$ and the sequence of $h_i, h^i$ is just $\{b, c\}$. Therefore, the former sequence together with the latter (read in reverse) comprise an alternating cycle. This final contradiction proves the claim.




\end{proof}
Now let $A$ be a finite presidium with no alternating cycles; we wish to show that $A$ embeds into the standard $\mathrm{TP}$-structure. First, suppose that $A$ has a convivium which is a singleton, $\{a\}$. Then $A \sqcup \{b\}$ with $b$ of sort $P$, where the instances of $R, I$ are exactly those of $A$ together with $I(b, a)$, will be a presidium with no alternating cycles. It follows from this that we may assume that each convivium of $A$ has size at least $2$. Similarly, we may assume that each convivium of $A$ has size at least $3$, even at least $4$.

Moreover, by the previous claim, we may now assume the following: if $a, b$ are distinct elements of $P(A)$ belonging to a common convivium, and $c, d$ are distinct elements of $P(A)$ belonging to a different common convivium, and if there is $e \in O(A) $ such that $R(e, b)$ and $R(e, c)$, then either there is $f \in O(A)$ (distinct from $e$) such that $R(f, a)$ and $R(f, c)$, or there is $f\in O(A)$ (again, distinct from $e$) such that $R(f, b)$ and $R(f, d)$.

With these assumptions in place, let $C_i$ and $C_j$ be two distinct convivia of $A$, and suppose that there are some $a \in C_i$, $b \in C_j$ and  $c \in O(A)$ such that $R(c, a)$ and $R(c, b)$.  Since $A$ has no alternating cycles (particularly, none with $1$ or $2$ points of sort $O$), there can be no distinct $a_1, a_2 \in C_i$, $b_1, b_2 \in C_j$ such that $a_1$ and $b_1$ have a common $R$-neighbor, and $a_2$ and $b_2$ have a common $R$-neighbor. Therefore, by the above assumption, it is either the case that for $C_i=\{a_\ell\}_{\ell \in I}$ with $a_i$ distinct, there are $\{c_\ell\}_{\ell \in I}$ with $c_\ell \in O(A)$ distinct such that $R(c_\ell, a_\ell)$ and $R(c_\ell, b)$, or that for $C_j=\{b_\ell\}_{\ell \in I}$ with $b_\ell$ distinct, there are $\{c_\ell\}_{\ell \in I}$ with $c_\ell \in O(A)$ distinct such that $R(c_\ell, b_\ell)$ and $R(c_\ell, a)$. (That the $c_{\ell}$ are distinct follows, say, from soundness of $A$.)


For $a, b \in P(A)$ belonging to distinct convivia, we now define $a < b$ if, for $\{b_i\}_{i \in I}$ the convivium to which $b$ belongs with $b_i$ distinct, there are $\{c_i \}_{i \in I}$ with $c_\ell \in O(A)$ (which must be distinct, by soundness) such that $R(c_i, b_i)$ and $R(c_i, a)$.  By the previous paragraph we now know that, if $a$ and $b$ have a common $R$-neighbor (and thus, by soundness, must belong to different convivia), either $a < b$ or $b < a$. Moreover, since $A$ has no alternating cycles (with $1$ or $2$ $O$-points) and each convivium has size at least $2$, the relation $<$ is antisymmetric: it is impossible that $a < b$ and $b < a$. We will also need that $<$ has no $3$-cycles: there are no $a < b$, $b < c$, $c < a$. To see this, because the convivium to which $c$ belongs has size at least $2$, there is $c'$ with $I^{*}(c, c')$. Because $b < c$, $c'$ and $b$ will have a common $R$-neighbor, $d$. Since the convivium to which $b$ belongs has size at least $3$, and $a < b$, we may find another member $b'$ (so $I^{*}(b, b')$) for which $b'$ and $a$ have a common $R$-neighbor $e$, distinct from $d$. Finally, since the convivium to which $a$ belongs has size at least $4$, and $c < a$, we may find another member $a'$ (so $I^{*}(a, a')$) for which $a'$ and $c$ have a common $R$-neighbor, $f$, distinct from $d$ and $e$. Then $b, b', d, c, c', e, a, a', f$ form an alternating cycle, a contradiction.


Our next goal will be to show that we can assume that $<$ is a tree ordering: $<$ is transitive and for $b < a$, $c< a$, either $b < c$ or $c < b$. To this end we show the following two claims:


\begin{claim}
    \label{order construction 2} Suppose $a < b$, $b < c$, and there is not already some $f\in O(A)$ with $R(f, a)$ and $R(f, c)$. Then the presidium $A' : = A \sqcup \{f\}$, $f \in O(A')$, whose instances of $R, I$ are exactly those of $A$ together with $R(f, a)$ and $R(f, c)$, has no alternating cycles.
\end{claim}


\begin{proof}
(of claim) 

(We do not really need ``there is not already some $f\in O(A)$ with $R(f, a)$ and $R(f, c)$" here, but include it to simplify the proof.) Because $a < b$ there are $b' \in P(A)$ with $I^{*}(b', b)$ (so $b'$ is distinct from $b$, and $b'$ is distinct from $a$ and $c$ because $b'$ is in the same convivium as $b$), and distinct $d, d' \in O(A)$ with $R(d, a)$, $R(d', a)$, $R(d, b)$, $R(d', b')$. Because $b < c$, there is also $e \in O(A)$  such that $R(e, b)$ and $R(e, c)$; $e$ is distinct from $d $ and $d'$ by the assumption that there is not already some $f\in O(A)$ with $R(f, a)$ and $R(f, c)$. Now suppose the presidium $A' : = A \sqcup \{f\}$ described in the claim has an alternating cycle. Then since $A$ does not already have an alternating cycle, there is a sequence 

  $$a_1, a_2, a^{1}, a_{3}, a_{4}, a^{2}, \ldots a_{2k-1}, a_{2k}, a^{k}, a_{2k+1}, a_{2k+2}, \ldots a_{2n-1}, a_{2n}, a^{n}, a_{2n+1}, a_{2n+2} $$

for $n \geq 1$ consisting of distinct $a_i \in P(A)$, $a^i \in O(A)$ with $I^{*}(a_i, a_{i+1})$ for $i \leq 2n+2$ odd, $R(a^i, a_{2i})$ and $R(a^i, a_{2i+1})$ for $i \leq n$,  and $a_1 = a$ and $a_{2n+1} = c$. Now suppose that the intersection of this sequence with $\{b, b', d, d', e\}$ is nonempty. First, suppose in particular that the first member of this intersection in the sequence (ordered as per the above enumeration), $a^k$, is either $d$ or $d'$. Then the interval between $a_1$ and $a^k$ inclusive will comprise an alternating cycle in $A$, a contradiction. Now suppose the same intersection is nonempty and its first member in the sequence, $a_k$, is $b$. If $k$ is odd, then the interval between $a_1$ and $a_k$ inclusive together with $b'$ and $d'$ comprises an alternating cycle, a contradiction. If $k$ is even, then $I^{*}(a_{k-1}, a_{k})$ and $I^{*}(a_k, b')$ implies $a_{k-1}$ and $b'$ are in the same convivium, so $I^{*}(a_{k-1}, b')$. Then the interval between $a_{1}$ and $a_{k-1}$ inclusive together with $b'$ and $d'$ comprises an alternating cycle, also a contradiction. The case where this same intersection is nonempty, and its first member in the sequence is $b'$, is handled similarly, with $b$ and $d$ in the role of $b'$ and $d'$, again getting a contradiction. If the intersection is nonempty and its first member in the sequence, $a^{k}$, is $e$, then the interval between $a_1$ and $a^{k}$ inclusive together with, say, $b$, $b'$ and $d'$ will comprise an alternating cycle, yet again a contradiction. So the intersection between this sequence and $\{b, b', d, d', e\}$, so in particular with $\{b,b',  d, e\}$, is nonempty. Therefore,  $a_i, a^{i}$ together with $e$, $b$, $b'$ and $d$ will comprise an alternating cycle; this final contradiction gives us the claim.

    
\end{proof}


\begin{claim}
    \label{order construction 3} Suppose $a < b$, $c < b$ with $a, c$ distinct, and there is not already some $f\in O(A)$ with $R(f, a)$ and $R(f, c)$. Then the presidium $A' : = A \sqcup \{f\}$, $f \in O(A')$, whose instances of $R, I$ are exactly those of $A$ together with $R(f, a)$ and $R(f, c)$, has no alternating cycles.
\end{claim}

\begin{proof}
    (of claim)

    (Again, we do not really need ``there is not already some $f\in O(A)$ with $R(f, a)$ and $R(f, c)$" here, but include it to simplify the proof.) 

    There is $b'$ distinct from $b$ in the same convivium as $b$, as well as $d, d', e, e' \in O(A)$, such that $R(d, a)$, $R(d', a)$, $R(d, b)$, $R(d', b')$ and $R(e, c)$, $R(e', c)$, $R(e, b)$, $R(e', b')$; any of the $d, d'$ must be distinct from any of the $e, e'$ by the assumption that  there is not already some $f\in O(A)$ with $R(f, a)$ and $R(f, c)$, while $c$ must be distinct from $c'$ and $d$ from $d'$ by soundness. Suppose that the presidium $A' : = A \sqcup \{f\}$ described in the claim has an alternating cycle. As  in the previous two claims, there is a sequence 

  $$a_1, a_2, a^{1}, a_{3}, a_{4}, a^{2}, \ldots a_{2k-1}, a_{2k}, a^{k}, a_{2k+1}, a_{2k+2}, \ldots a_{2n-1}, a_{2n}, a^{n}, a_{2n+1}, a_{2n+2} $$

for $n \geq 1$ consisting of distinct $a_i \in P(A)$, $a^i \in O(A)$ with $I^{*}(a_i, a_{i+1})$ for $i \leq 2n+2$ odd, $R(a^i, a_{2i})$ and $R(a^i, a_{2i+1})$ for $i \leq n$,  and $a_1 = a$ and $a_{2n+1} = c$. Suppose that the intersection of this sequence with $\{b, b', d, d', e, e'\}$ is nonempty. First, suppose in particular that the first member of this intersection in the sequence (once again, ordered according to the above enumeration), $a^{k}$ is one of $d, d', e, e'$. If $a^{k} $ is $d$ or $d'$, then the interval from $a_1$ to $a^k$ inclusive in this sequence will comprise an alternating cycle, a contradiction. If $a^{k}$ is $e$, the interval from $a_1$ to $a^k$ inclusive together with $b, b', d'$ will comprise an alternating cycle, a contradiction, while if $a^{k}$ is $e'$, the interval from $a_1$ to $a^k$ inclusive together with $b', b, d$ will comprise an alternating cycle, also a contradiction. Now suppose this first member, $a_k$, is $b$; if $k$ is odd then the interval from $a_1$ to $a_k$ inclusive together with $b'$ and $d'$ will comprise an alternating cycle, a contradiction. If $k$ is even then $I^{*}(a_{k-1}, a_{k})$ and $I^{*}(a_{k}, b')$ will imply $a_{k-1}$ and $b'$ are in the same convivium, so $I^{*}(a_{k-1}, b')$. Then the interval from $a_1$ to $a_{k-1}$ inclusive together with $b'$ and  $d'$ will comprise an alternating cycle, again a contradiction. (This is exactly as in the case where $a_{k}$ is equal to $b$ in the previous claim.) The case where this first member is $b'$ is handled similarly, with $b$ and $d$ in the role of $b'$ and $d'$, again getting a contradiction. So the intersection of the sequence with $\{b, b', d, d', e, e'\}$, and particularly with, say, $\{b, b', d', e\}$, must be empty. The sequence of $a_i, a^i$ together with $e$, $b$, $b'$, and $d'$ must comprise an alternating cycle, a contradiction that proves our claim. (So we did not need, say, $e'$ here at all, but include it anyway for reasons of symmetry.)
\end{proof}

Then, applying our three claims in this proof, we may assume that $A$ satisfies the assumptions from the paragraph after the proof of Claim \ref{order construction 1}, and also that any distinct $a, b$, such that there is a $c$ with $a < c < b$ or $a < c$, $b < c$, must have a common $R$-neighbor. (Assuming that $A$, and thus the number of pairs of points of $P(A)$, is finite, we only need to add finitely many points of $O(A)$, so $A$ stays finite). So, by the discussion preceding Claim \ref{order construction 2}, $<$ is, in addition to being antisymmetric, transitive and a tree order. (In particular, the absence of $3$-cycles for $<$ is used to prove transitivity: if $a< b$, $b< c$, then either $a < c$ or $c < a$, because $a$, $c$ must have a common $R$-neighbor. But $c <a$ is ruled out by the absence of $3$-cycles.)


Now define an injective map $\iota$ of $P(A)$ into the sort-$P$ points $\{b_{\sigma}\}_{\sigma \in \eta}$ of the standard $\mathrm{TP}$-structure, inductively on the $<$-height of an element of $P(A)$, as follows. Note that by definition of $<$, all elements of the same convivia have the same $<$-downsets, so have the same $<$-height. First, let $C_0, \ldots, C_\ell$ enumerate the convivia of $A$ all of whose elements have height $0$.  Then define $\iota(b)$, for $b$ the $k$-th element of $C_{i}$ in the linear order on $C_i$ defined by $I$, to be equal to $b_{\langle ik \rangle}$. Now suppose $\iota$ has been defined on elements of $P(A)$ of $<$-height at most $n-1$. For each element $b'$ of $<$-height $n -1$, let $C^{b'}_{0}, \ldots, C^{b'}_{\ell^{b'}}$ be a fixed enumeration of the convivia whose elements are all immediate $<$-successors of $b$. Let $b$ be an element of $P(A)$ of $<$-height $n > 0$; we define $\iota$ on $b$. Since $<$ is a tree order, there is a unique element $b'$ of $P(A)$ of height $n-1$ with $b' < b$.  Suppose we have already defined $\iota(b') = b_{\eta}$. Then define $\iota(b) := b_{\eta \smallfrown \langle ik \rangle}$, where $b$ is in $C^{b'}_i$ and is the $k$-th element of that convivium in the linear order on it defined by $I$. Having defined this map, we see it has the following two features: first, $I(b_1, b_2)$ implies $I(\iota(b_1), \iota(b_2))$ for $b_1, b_2 \in P(A)$, and second, if $b_1 < b_2$ for $b_1, b_2 \in P(A)$, and $\iota(b_1)=b_{\eta_1}$ while $\iota(b_2) = b_{\eta_2}$, then $\eta_{1} \lhd \eta_2$.


It remains to extend this map $\iota$ on $P(A)$ to an embedding of all of $A$ into the standard $\mathrm{TP}$-structure. We can find a family of distinct points $a^i_{S}$ of sort $O$ of the standard $\mathrm{TP}$-structure, where $i < \omega$ and $S$ ranges over those finite subsets of $\omega^{< \omega}$ linearly ordered by $\lhd$, such that $R(a_{S}^{i}, b_{\eta})$ for $\eta \in S$. Now, for any $A_{0} \subset P(A)$, let $\{a_{A_{0}}^{i}\}^{n_{A_{0}}}_{i =0}$ be an enumeration, with $a^{i}_{A_{0}} \neq a_{A_{0}}^{i'}$ for $i \neq i'$, of the set of points in $O(A)$ consisting of those $a$ for which $A_0$ is equal to the set of all $b \in P(A)$ with $R(a, b)$. By definition, for distinct $A_{0}, A'_{0}$, the set $\{a_{A_{0}}^{i}\}^{n_{A_{0}}}_{i =0}$ and $\{a_{A'_{0}}^{i}\}^{n_{A'_{0}}}_{i =0}$ are disjoint.  Then $\{a_{A_{0}}^{i}\}^{n_{A_{0}}}_{i =0}$ can only be nonempty if $A_{0}$ is linearly ordered by $<$, by the discussion before Claim \ref{order construction 2}.  For any linearly ordered $A_{0} \subset P(A)$, let $S_{A_0}$ be the set of $\eta \in \omega^{<\omega}$ such that there is $b \in A_{0}$ with $\iota(b) = b_{\eta}$; by the second condition satisfied by the map $\iota$ on $P(A)$, $S_{A_0}$ is linearly ordered by $\lhd$. So we define $\iota(a^i_{A_{0}}):=a_{S_{A_{0}}}^{i}$ for $A_0$ any subset of $P(A)$. Then $\iota$, an injective map defined on all of $A$, satisfies the condition that $R(a, b)$ implies $R(\iota(a), \iota(b))$ for $a \in O(A)$, $b\in P(A)$. With the condition that $\iota$ was noted to satisfy before that $I(b_1, b_2)$ implies $I(\iota(b_1), \iota(b_2))$, we have shown that $\iota$ is a weak embedding of $\mathcal{L}_{\mathrm{std}}$-structures, completing the proof of the lemma.

\end{proof}

\end{document}
